\documentclass[11pt]{article}
\usepackage[margin=1in]{geometry}
\usepackage[T1]{fontenc}
\usepackage{lmodern}
\usepackage{hyperref}
\usepackage{booktabs}
\hypersetup{colorlinks=true, linkcolor=blue, urlcolor=blue, citecolor=blue}

\title{Nature, Security \& Defense:\\A Survey of a Unified Research Programme}
\author{Sefunmi Ashiru\\Architecture of Homo Sapiens LLC}
\date{November 2026\\Preprint (draft)}

\begin{document}
\maketitle

\begin{abstract}
This paper surveys nature, security, and defense as a single research
programme. It brings together five strands that are usually treated
separately: nature-inspired defense engineering (biomimicry, bat
echolocation, swarms, and immune-inspired security); acoustic defense
and sonar; environmental and ecological security; and verification and
trust in defense systems (provenance, attestation, and hash anchoring).
The strands converge on three shared questions: resilience without a
centre, sensing and verification under noise and deception, and nature
as both model and stakes. The survey is grounded in cited sources ---
journal articles, conference papers, and institutional assessments ---
and grades the evidence by strand rather than blending it. Sources the
author read in full are marked \emph{verified live}; sources read only
as search-result text are marked \emph{index} and must be re-checked
before final citation. A final section records concrete improvements
the cited research suggests for the development of a software platform
whose defense and security layer relies on signed commits, provenance
and SBOMs, and blockchain-anchored publish receipts, labelling each
transfer \emph{Direct}, \emph{Analogy}, or \emph{Speculative}.
\end{abstract}

\section{Introduction}

Defense research has long borrowed from nature. Camouflage imitates
animal coloration; sonar imitates bat echolocation; swarm tactics
imitate flocking and schooling; intrusion detection imitates the
immune system. At the same time, nature itself has become a security
domain: climate change, resource scarcity, and ecosystem loss are now
treated by NATO, the IPCC, and UNEP as security questions, not only
environmental ones. And the systems that defend anything --- a drone, a
supply chain, a software release --- must now \emph{prove} what
happened, through provenance, attestation, and tamper-evident records.

This paper treats those three facts as one programme rather than three
fields. Section 2 covers nature-inspired defense engineering;
Section 3 covers acoustic defense and sonar; Section 4 covers
environmental and ecological security; Section 5 covers verification
and trust; Section 6 states what a unified programme would study, its
open questions, and where the fields disagree; Section 7 records
implications for software development; Section 8 records limitations
and what could not be verified.

Two cautions run throughout. First, no source was found that studies a
unified ``nature, security and defense'' programme as such; the
synthesis in Section 6 is the author's, built from the strand sources.
Second, evidence quality is uneven across strands: Sections 3 and 5
rest on peer-reviewed, replicated results; Section 2 rests heavily on
reviews, programme reports, and preprints; Section 4 rests on
institutional assessments rather than experiments.

\section{Nature-Inspired Defense and Security}
\label{sec:biomimicry}

\subsection{Key findings}

\textbf{Minimal bat sonar works, at short range and low speed.}
Bats navigate with one emitter and two receivers using broadband
frequency-modulated (FM) pulses. Yamada et al.\ reproduced this with a
6.6~g thermophone plus two MEMS microphones (343~g total system),
achieving about 17~mm ranging resolution by cross-correlation, and
low-speed two-dimensional obstacle avoidance during autonomous flight
even when microphone signals were dominated by rotor noise
(about 94~dB SPL at 0.7~m) \cite{yamada2026}. This is a verified-live
source.

\textbf{Low-power ultrasound beats vision where vision fails, and
fails where vision succeeds.} BatDeck (ETH Zurich) achieved
ego-velocity mean squared error of 0.019~m/s and 8-minute / 136~m
exploration among transparent and reflective obstacles, using sensors
drawing under 1~mW. Its own comparison table shows cameras, laser
time-of-flight, and ultrasound each fail differently, so robustness
comes from complementing them, not replacing one with another
\cite{batdeck2024}. This is a verified-live source.

\textbf{Bio-inspired materials are a parallel, more mature strand.}
Nacre, mantis-shrimp dactyl clubs, cephalopod skin, and armadillo
shells inspire armour, self-healing, energy dissipation, and adaptive
camouflage; scalability, cost, and reliability under harsh conditions
remain the stated bottlenecks \cite{igi2025,armadillo2024}. These are
index sources.

\textbf{Mimicry is also concealment, deception, and imitation.}
Matthews and Matthews (2024) trace military camouflage, deception with
fakes and dummies, and imitation (animals, bioengineering, avian,
insect, and mammal sonar) toward ``robots, swarms, and avatars'' driven
by AI \cite{matthews2024}. Index source.

\textbf{Swarm tactics have been demonstrated at scale by DARPA.}
The OFFSET programme envisioned swarms of up to 250 collaborative
autonomous air and ground systems for small units in urban areas, with
a Swarm Tactics Exchange for composing collective behaviours,
algorithms, and tactics \cite{darpa2020}. Index source.

\textbf{Immune-inspired cybersecurity is a long literature with open
scalability questions.} The human immune system is invoked for robust,
self-organised, distributed protection; artificial immune systems (AIS)
use negative and clonal selection, self/non-self discrimination, and
immune memory for intrusion detection (notably the LISYS system).
Reviews conclude substantial potential but call for work on scalability
and on clonal selection and immune memory \cite{kim2007,bala2025}.
Index sources.

\subsection{Implications for development}
\label{sec:biomimicry-dev}

\textbf{Complementary, fail-differently controls --- Analogy, grounded
in \cite{batdeck2024}.} BatDeck's comparison shows no single sensor is
robust; the system is robust because sensors fail differently. A
software platform already has this shape (signed commits, SBOM,
SAST/secret scanning, anchored receipts). The sourced improvement is to
\emph{test the combination for correlated failure}: a single compromised
CI runner or a single signing key that defeats several checks at once is
the equivalent of a sensor suite sharing one blind spot. Do not add a
fifth check that shares the same trust root as the existing four.

\textbf{No single point of failure / quorum --- Analogy, grounded in
\cite{darpa2020} and \cite{swarmmesh2026}.} OFFSET's swarm has no
central controller whose loss ends the mission; the Swarm-Mesh preprint
proposes quorum thresholds so no single agent can trigger containment or
be compromised undetected. Transferred to a platform: require more than
one independent approval or signature for a publish (for example, author
signature plus CI builder attestation), and make publish fail if either
is missing. Speculative as a software claim: no source in this strand
tests quorum approval for software releases.

\textbf{Self/non-self anomaly detection --- Analogy, grounded in
\cite{kim2007}.} AIS distinguishes ``self'' from ``non-self'' and
adapts. For a platform this supports an allowlist model already implicit
in SBOMs: a release whose dependency set differs from the recorded SBOM
should be treated as ``non-self'' and fail closed, rather than merely
warned about. The AIS literature itself warns scalability is unresolved,
so this should be a deterministic SBOM-diff check, not a learned anomaly
detector, unless separately validated.

\textbf{Do not over-claim bio-inspiration.} A bio-inspired blockchain
paper discussed in Section \ref{sec:verification} claims 95.2\%
anomaly detection in simulation; that is a simulation claim, not a
field result, and should not be cited as evidence a platform control
works.

\section{Acoustic Defense and Sonar}
\label{sec:acoustic}

\subsection{Key findings}

\textbf{Saranga (WPI, \emph{Science Robotics} 2026) is the strongest
result.} A 16~cm, roughly \$400 palm-sized quadrotor using only 1.2~mW
of sensing power and onboard computation navigated dense fog, darkness,
and snow among thin and transparent obstacles, using a dual sonar array.
It overcame a peak signal-to-noise ratio of $-4.9$~dB with two combined
fixes: (a) a physical acoustic shield (inspired by bumblebee-bat ear
cartilage) blocking propeller noise, and (b) a deep neural network
denoising weak echoes using the long horizon of ultrasound echoes.
Training generalised via a synthetic-data pipeline plus limited real
noise data. Across 180 test flights success was 72--100\%; it struggled
with thin metal poles and slender branches that weakly reflect
\cite{saranga2026}. Verified live (arXiv abstract).

\textbf{Thermophone FM sonar (Yamada et al.\ 2026) shows the same noise
problem solved differently:} broadband FM pulses plus
cross-correlation pulse-compression gain plus thresholding, detecting a
12~cm pole while hovering under dominant rotor noise
\cite{yamada2026}. Verified live.

\textbf{Synthetic echolocation (Michigan, Army-funded) trains on
simulation alone.} Patil et al.\ (2025, \emph{Journal of Sound and
Vibration} 119229) and Kwon et al.\ (2025, \emph{IEEE Access}) trained
ensembles of convolutional neural networks --- each network for one
shape --- entirely on numerically simulated echoes, and classified real
measured echoes, distinguishing spheres from cylinders despite similar
curved frontal geometry; low-symmetry cubes head-on were the failure
case. Funded by the US Army Research Office and DEVCOM Ground Vehicle
Systems Center \cite{michigan2025}. Index source (University of
Michigan ARC project listing); the papers themselves were not opened.

\textbf{Hierarchical reinforcement learning is the WPI navigation
stack.} The NSF-funded Sanket project (from September 2025, \$704,908,
3 years) uses physics-informed deep learning to interpret ultrasound and
a hierarchical reinforcement-learning stack to move toward goals while
avoiding obstacles, fused with IMU and other sensors, for robots under
100~mm / 100~g, with real-world deployment aimed at 3--5 years
\cite{wpi2025}. Index source.

\textbf{Active noise cancellation for naval sonar: no strong
peer-reviewed source was found.} Patent and trade material describes
active cancellation (anti-phase counter-waves, reference hydrophones,
adaptive algorithms) for hull, propeller, and torpedo-seeker noise, and
acoustic cloaking and metamaterials as laboratory proof-of-concept with
narrow bandwidth. This is a \emph{gap}, not a finding --- the patent
analysis must not be cited as peer-reviewed evidence. The closest
peer-reviewed analogue is Saranga's physical shield plus learned
denoising, above.

\subsection{Implications for development}

\textbf{Two-layer noise handling --- Analogy, grounded in
\cite{saranga2026}.} Saranga succeeds because it combines a
\emph{deterministic physical} fix (shield) with a \emph{statistical}
fix (learned denoising); neither alone sufficed at $-4.9$~dB SNR. For CI
scanning, the analogue is: deterministic allowlists and baselines that
suppress known-benign findings (the ``shield'') \emph{plus} a second,
independent check for what remains --- not two copies of the same
scanner. Speculative: no source tests this transfer to SAST/secret
scanning; treat as a design hypothesis, not evidence.

\textbf{Synthetic-data testing of the verifier --- Analogy, grounded in
\cite{saranga2026,michigan2025}.} Both Saranga and the Michigan work
generalise by training and testing on synthetic echoes plus limited real
data. For a platform: generate synthetic tampered artifacts (mutated
SBOM, re-signed commit, hash-mismatch receipt, swapped dependency) and
require the verification path to catch them --- that is, test the
verifier, not just the happy path. This is a direct methodological
transfer, but the \emph{effectiveness} of synthetic testing for software
verification is not established by these papers.

\textbf{Hierarchical checks (cheap local to expensive anchored) ---
Speculative, loosely grounded in \cite{wpi2025}.} WPI's hierarchical
reinforcement learning does fast local avoidance before higher-level
goal navigation. Analogously: run cheap local verification (recompute
hash, check signature) on every publish, and reserve expensive anchored
verification (on-chain receipt lookup, full SBOM diff) for
releases and promotion gates. Flagged as speculative: the source does
not discuss verification hierarchies.

\textbf{What this strand does \emph{not} support:} active noise
cancellation as a software pattern. Do not claim ``active noise
cancellation'' as a sourced platform technique; the sourced techniques
are physical shielding, cross-correlation and thresholding, and learned
denoising.

\section{Environmental and Ecological Security}
\label{sec:ecological}

\subsection{Key findings}

\textbf{Climate change is a threat multiplier, not a direct cause of
war --- the consensus position.} IPCC AR6 WGII (2022) finds climate
extremes are not directly causing new conflicts in large statistical
studies; they prolong and increase the severity and frequency of ongoing
violent conflicts, with a small but persistent signal, stronger where
livelihoods depend on climate-sensitive agriculture and livestock and
governance is weak. Higher warming increases conflict risk via food and
livelihood insecurity \cite{ipcc2022,newsecuritybeat}. Index sources.

\textbf{NATO has institutionalised this.} Its 2021 Climate Change and
Security Action Plan calls climate change ``one of the defining
challenges of our times'' and ``a threat multiplier'', and commits to
four pillars: an annual Climate Change and Security Impact Assessment
(awareness); adaptation into resilience, defence planning, capability
delivery, standards, procurement, and training; mitigation via a
greenhouse-gas-emissions mapping methodology for military activities;
and outreach \cite{nato2021}. Verified live. NATO's topic page adds
environmental security (extreme weather, sea-level rise, resource
depletion, land degradation, pollution leading to humanitarian disasters,
regional tensions, and violence), the Science for Peace and Security
programme (since 1969), Trust Funds for demilitarisation and demining,
and a Code of Conduct for active sonar to protect marine mammals,
developed after CMRE research on sonar effects \cite{nato-topic}.
Verified live.

\textbf{Natural resources are central to conflict and to
peacebuilding.} UNEP (2009): 40\% of intrastate conflicts in the past
60 years were strongly linked to natural resources; such conflicts are
twice as likely to relapse within five years; fewer than a quarter of
peace agreements address natural-resource issues \cite{unep2009}. Index
source. UNEP (2024), from 40 case studies, finds nature-based solutions
(protect, conserve, restore, and sustainably use natural resources) are
``uniquely placed to address underlying conflict drivers'', while
conflict itself directly threatens such projects, and inequality,
economic incentives, and migration must be addressed \cite{unep2024}.
Index source.

\textbf{Protection of nature is itself a security task.} NATO's sonar
code of conduct is a concrete case of a defence capability being
constrained to protect a natural system; UNEP's ``Greening the Blue
Helmets'' work applies the same logic to peacekeeping's own
environmental footprint \cite{nato-topic,unep2009}.

\subsection{Implications for development}

\textbf{Threat-multiplier risk framing --- Analogy, grounded in
\cite{ipcc2022,nato2021}.} IPCC and NATO insist a stressor rarely
\emph{causes} failure alone; it amplifies existing fragility. For a
platform: treat a vulnerable dependency, a leaked secret, or an unsigned
commit not as an isolated CVE to patch but as a multiplier on existing
weaknesses (for example, a vulnerable crate matters far more if the same
build can publish without a second signature). Prioritise fixes by
\emph{interaction} with existing fragility, not by severity score alone.
Analogy, not a sourced software method.

\textbf{Annual / per-release impact assessment --- Analogy, grounded in
\cite{nato2021}.} NATO conducts an \emph{annual} Climate Change and
Security Impact Assessment and integrates it into risk and resilience
assessments and procurement. Analogously: a scheduled, recorded
security-impact assessment per release (SBOM diff plus new dependencies
plus changed attestations), stored alongside the anchored receipt, rather
than ad-hoc review. This is a process transfer; no source tests it for
software.

\textbf{Address underlying drivers, not only symptoms --- Analogy,
grounded in \cite{unep2024}.} UNEP's lesson is that nature-based
solutions work when they address underlying conflict drivers
(inequality, incentives, livelihoods). For a platform: the analogue of
``underlying drivers'' is dependency hygiene and build hygiene (pinning,
minimal dependencies, reproducible builds) rather than adding another
scanner on top of an undisciplined supply chain. Speculative as a causal
claim.

\textbf{Minimise collateral / data minimisation --- Direct-ish, grounded
in \cite{nato-topic} and Section \ref{sec:verification}.} NATO
constrained active sonar to protect marine mammals; the verification
literature (Section \ref{sec:verification}) likewise holds that only a
hash should leave the machine. For a platform: publish only
digests and commitments, keep SBOMs and provenance evidence locally, and
treat any collection of source or user data in a receipt as collateral
to be avoided. This is supported by the anchoring sources in
Section \ref{sec:verification}, not by the ecological sources alone.

\section{Verification and Trust in Defense Systems}
\label{sec:verification}

\subsection{Key findings}

\textbf{SLSA gives the normative build-provenance ladder.}
Specification v1.2, Build track: L1 provenance exists; L2 signed and
authentic provenance from a hosted build service; L3 unforgeable
provenance plus isolated and ephemeral builds. The producer must choose
a capable platform, follow a consistent build, and distribute
provenance. Provenance must identify the output by cryptographic digest.
Critically: at L2 provenance \emph{must} be generated by the control
plane, not by a tenant or build step, and the signing key must be
accessible only to the build service --- not to user-defined build
steps; at L3 every field is generated and verified in the trusted
control plane, external parameters are fully enumerated, and
cross-build and cache-poisoning influence is prohibited. Transparency
logs or timestamping are recommended for key-compromise detection
\cite{slsa}. Verified live.

\textbf{Hash anchoring proves a narrow thing.} Committing a hash (or a
Merkle root of many hashes) to a public chain proves the exact bytes
existed at or before the block time and have not changed since;
verification needs only the document, a small proof file, and public
chain data. It does \textbf{not} prove authorship, accuracy,
exclusivity (ten contradictory drafts can be anchored), legal execution,
or precise time (hour-level for Bitcoin) \cite{filigrane2026}. Verified
live. Haber and Stornetta (1991) hash-chaining predates blockchains;
OpenTimestamps aggregates unbounded digests into one Bitcoin
transaction, so individual timestamps are effectively free.

\textbf{The correct pattern is: anchor roots, keep proofs, verify on a
schedule, combine instruments.} Filigrane's field notes: anchor Merkle
roots not raw hashes (privacy amplification, batching); automate
\texttt{.ots} proof upgrade after confirmation; store proofs with the
documents; verify on a schedule (re-hash, walk the Merkle path, check
the block header); plan for algorithm ageing (SHA-1 was retired);
combine signature (origin) plus qualified TSA token (legal time) plus
chain anchor (independent longevity) \cite{filigrane2026}. Verified live.

\textbf{A worked defence-supply-chain analogue exists (vendor white
paper).} Sovereign Proof's Sentinel gives every serialized asset a
signed, hash-chained custody biography (Ed25519, SHA-256, genesis
value), a fleet Merkle root, and a single 32-byte commitment anchored to
Bitcoin via OpenTimestamps, while all custody evidence stays inside the
operator's facility. Its principles: determinism over judgment (an LLM
narrates evidence, never computes a score); evidence stays home;
tamper-evidence over tamper-prevention; fail closed (a missing
attestation library is an error, never a silent downgrade to software);
no token. Mutating a record breaks that record while the rest verifies;
the fleet root then diverges from the anchor. Silicon binding uses TPM
2.0 quotes / DICE / SPDM over a challenge nonce; a quote from the wrong
key is a fatal mismatch. It maps evidence to DFARS 252.246-7007/-7008
and NDAA FY2023 \S5949 \cite{sovereign2026}. Verified live; vendor white
paper, not peer-reviewed.

\textbf{Hardware attestation is the same pattern at the firmware
layer.} NIST SP 800-193 (2018) organises platform firmware resiliency
as Protection / Detection / Recovery. Secured boot only lets
authenticated firmware execute, starting from a hardware root of trust;
measured boot records hashes into TPM PCRs so a remote verifier can
compare them to golden measurements (remote attestation)
\cite{nist800193}. Index source (secondary summaries); the SLSA
specification should be the primary citation for provenance.

\textbf{Sigstore / in-toto / Rekor complete the software stack.}
Sigstore issues short-lived (about 10 minute) certificates bound to an
OIDC identity (Fulcio), signs with Cosign, and logs every signing event
in Rekor, an append-only Merkle transparency log. in-toto defines a
\emph{layout} (authorised steps, keys, and rules) plus signed
\emph{links} recording hashes of each step's inputs and outputs, so
verification checks an unbroken signed chain from named source to
artifact. SBOM formats are SPDX (ISO/IEC 5962) and CycloneDX (OWASP)
\cite{sigstore-guide}. Index, secondary sources; verify against
sigstore.dev and in-toto.io before final use.

\textbf{Passkeys are FIDO2/WebAuthn credentials, not passwords.}
A passkey is a WebAuthn credential: the authenticator (a phone,
laptop, or hardware security key) generates a per-service
public-private key pair, the private key never leaves the
authenticator, and the server stores only the public key.
Authentication is a challenge-response ceremony bound to the relying
party's origin, so a credential registered for one domain cannot be
used on a phishing copy of that domain. That makes passkeys
phishing-resistant by construction, unlike passwords and one-time
codes, which can be relayed \cite{webauthn3,nist80063b4}. WebAuthn
Level 3 became a W3C Recommendation on 25 August 2026, superseding
Level 2 (2021); it formalises passkeys at scale, including
backup-eligible / backup-state flags that distinguish synced from
device-bound credentials \cite{webauthn3}. NIST SP 800-63B-4 (July
2025) admits syncable authenticators at AAL2 but not at AAL3, which
requires a non-exportable key --- the formal basis for preferring
device-bound passkeys for privileged accounts \cite{nist80063b4}.
Verified live. The distinction that matters for a
platform is scope: passkeys authenticate a \emph{human} to a service
through a browser or operating system. They do not replace a network
address (an RPC URL or a contract address is a locator, not a secret,
and is public by design), and they do not by themselves replace a
\emph{service} key held by a server process that must sign when no
human is present.

\textbf{Tamper-evidence is explicitly not tamper-prevention, and cannot
defeat a fully compromised signer.} STAMP (molecular provenance
preprint) states this plainly: it does not claim to defeat jailbroken
synthesizers --- that is irreducible --- it is a cost imposer and
evidence generator that converts attacks into forensically suspicious
artifacts or supply-chain-visible events \cite{stamp}. Index source;
preprint, authors not captured.

\subsection{Implications for development --- mostly Direct}

These map onto a platform's existing controls (signed commits, verified
author identity, provenance plus SBOMs, SAST/secret scanning,
blockchain-anchored receipts) and are the most actionable part of this
survey.

\textbf{Generate provenance in the control plane, not in the build ---
Direct, \cite{slsa}.} SLSA L2/L3 require the provenance generator to be
the build platform itself, with the signing key inaccessible to
user-defined build steps. \emph{Improvement:} the platform's
provenance and SBOM attestation must be emitted by the CI control plane
(for example, GitHub Actions or the release service), never by
application code or a build script that a compromised dependency could
influence. If a build step can write its own provenance, the control is
L1 at best.

\textbf{Aim explicitly at SLSA Build L2, then L3 --- Direct,
\cite{slsa}.} L2 is a hosted build plus signed provenance; L3 is
isolated and ephemeral builds plus unforgeable provenance plus no
cache-poisoning or cross-build influence. \emph{Improvement:} state the
target level in the platform's security documentation, and treat
``builds on a developer laptop'' and ``shared, non-ephemeral runners
with a persistent cache'' as explicit L2/L3 failures.

\textbf{Anchor Merkle roots, not raw artifact hashes --- Direct,
\cite{filigrane2026,sovereign2026}.} Both Filigrane and Sovereign Proof
anchor a single root (fleet or root commitment) rather than one
transaction per document. \emph{Improvement:} anchor one Merkle root
per release covering the artifact hash plus SBOM hash plus provenance
hash plus commit hash, and keep the per-artifact proof (Merkle path)
with the release. This gives privacy amplification (an observer cannot
count documents), negligible marginal cost, and a single batching point.

\textbf{Keep the proof, upgrade it, verify on a schedule --- Direct,
\cite{filigrane2026}.} An OpenTimestamps-style receipt starts
incomplete and must be upgraded after confirmation; a lost proof is a
lost timestamp even though the anchor is permanent. \emph{Improvement:}
automate proof upgrade and retention with the artifact, and run
scheduled re-verification (recompute hash, walk Merkle path, check
block header) --- evidence nobody has ever verified tends to fail the
first time it matters \cite{filigrane2026}.

\textbf{Do not let the receipt claim what it cannot prove --- Direct,
\cite{filigrane2026,stamp}.} An anchor proves existence plus integrity
at a coarse time only. It does not prove authorship, accuracy,
exclusivity, or precise time. \emph{Improvement:} the platform's
``verified by recomputing the hash'' check must be presented as
existence and integrity only. Authorship needs the signed commit plus
verified author identity (separate signature); uniqueness and sequence
need an append-only log (cf.\ Rekor, \cite{sigstore-guide}), not just an
anchor; precise time needs an RFC 3161 TSA token if it matters. A false
statement anchored on-chain is a false statement with a verifiable date
\cite{filigrane2026}.

\textbf{Fail closed, never silently downgrade --- Direct,
\cite{sovereign2026,nist800193}.} Sovereign Proof: a missing
attestation library is an error, never a silent downgrade to software; a
key that cannot decrypt the keystore refuses to mint keys. NIST SP
800-193: detection plus recovery, not just protection.
\emph{Improvement:} if signature verification, SBOM generation, or
receipt verification is unavailable, the publish or deploy must fail ---
not proceed unsigned. Log the failure as a finding, not a warning.

\textbf{Deterministic scoring; evidence stays home --- Direct,
\cite{sovereign2026}.} Trust must be computed by published,
reproducible rules, never by a model's opinion; only a 32-byte
commitment leaves the facility. \emph{Improvement:} the platform's
verification verdict should be a deterministic function of signatures
plus hashes plus log inclusion (reproducible by any verifier from public
data), and SBOM and provenance content should not be published on-chain
--- only its digest or root.

\textbf{Bind identity to silicon or key, with a challenge nonce ---
Direct, \cite{sovereign2026,nist800193}.} TPM quotes over a fresh
nonce, verified against the key the device was born with.
\emph{Improvement (for the mobile and server components):} where a
device or build runner holds a key, verify a fresh challenge and
response at publish time rather than accepting a stored ``I am
authorised'' claim; a quote from the wrong key is a fatal mismatch, not
a warning.

\textbf{Use passkeys for human signing; keep a dedicated service key
for the gateway --- Direct for the passkey standard, Analogy for the
hybrid, \cite{webauthn3,nist80063b4,sovereign2026,nist800193}.}
\emph{Improvement:} human actions on the platform --- release
approval, quest completion, and any publish that a person authorises
--- should be signed with a passkey (device-bound for privileged
accounts, synced passkeys only for general use, with a registered
second authenticator as the recovery path for device loss), so no
shared secret is ever transmitted or stored server-side. The chain
gateway, by contrast, is a server process that must anchor a release
when no human is at the keyboard; it cannot complete a WebAuthn
ceremony, so it still needs a dedicated service key delivered through
hosting secrets, never in code or in CI. That key must be scoped to
the one thing it may do --- call \texttt{anchorArtifact} on the World
Registry Contract --- hold no assets, be rotatable, and have every
anchor it makes visible in the transparency log and the indexer, so a
compromised gateway key produces a visible bogus anchor rather than a
silent one. Passkeys and gateway secrets are therefore complementary,
not alternatives: passkeys remove the human shared secret, the service
key remains the machine shared secret, and the two must never share a
trust root. Where a publish needs both a person and the gateway,
require both signatures (passkey approval plus gateway anchor) so
neither alone can publish --- the quorum pattern from Section
\ref{sec:biomimicry}, applied to signing.

\textbf{State honest limitations --- Direct,
\cite{sovereign2026,stamp}.} Sovereign Proof and STAMP both publish a
limitations section (hardware backends are roadmap, not shipped;
jailbroken signers are irreducible). \emph{Improvement:} the platform
should document what its receipts do \emph{not} defeat --- for example,
a compromised signing key, a compromised CI control plane, or a
malicious authorised author --- and what would make such an attack
visible (transparency log, anchor divergence).

\section{How the Strands Connect}
\label{sec:connection}

\subsection{What a unified programme would study}

The strands converge on three shared questions, each visible in the
sources above.

\textbf{1. Resilience without a centre.} How do systems survive the
loss of components? Bats navigate with a minimal sensor set; swarms
(OFFSET) and quorum systems (Swarm-Mesh) survive node loss; immune
systems are distributed and self-organised; SLSA, Rekor, and
hash-anchoring survive the loss of any single vendor because
verification needs only public data. The programme's core hypothesis ---
\emph{redundancy plus local rules plus verifiable evidence beats a
single hardened centre} --- is the common thread, but it is a
hypothesis, not a proven theorem across domains.

\textbf{2. Sensing and verification under noise and deception.}
Ultrasound at $-4.9$~dB SNR, climate signals inside noisy conflict
statistics, and tampered provenance inside a noisy supply chain are the
same epistemic problem: extracting a trustworthy signal when the channel
is degraded and the adversary may be active. Saranga's two-layer answer
(deterministic shield plus statistical denoising) and SLSA's answer
(control-plane provenance plus transparency log) are domain-specific
instances of the same pattern.

\textbf{3. Nature as both model and stakes.} Nature supplies the design
(echolocation, swarm, immunity) \emph{and} is the thing being defended
(marine mammals versus sonar, ecosystems versus climate-driven conflict,
resources versus war). NATO's sonar code of conduct is the clearest
institutional example of holding both roles at once.

\subsection{Open questions}

\begin{itemize}
\item \textbf{What is the right level of bio-fidelity?} Yamada et al.\
note a ``fundamental gap between biological and artificial airborne
sonar'' --- bats adjust pulse intensity, time--frequency structure, and
beam width dynamically; current drones do not, and fly at 0.4~m/s. Is
closer imitation the path, or is the bat a proof that minimal
architectures suffice?
\item \textbf{Can immune and swarm ideas scale to real software
systems?} AIS reviews (2007, 2025) still list scalability, clonal
selection, and immune memory as open; the bio-inspired blockchain
paper's 95.2\% / 15.2~ms figures are simulation claims. No source in
this survey demonstrates a fielded AIS or swarm-immune system at
production scale.
\item \textbf{Does anchoring scale as a social institution, not just a
cryptographic one?} OpenTimestamps can anchor billions of hashes
cheaply, but the Filigrane analysis shows the binding constraints are
proof custody, algorithm ageing, legal standing (eIDAS qualified
timestamps versus admissible-but-unprivileged anchors), and whether
anyone actually verifies on a schedule.
\item \textbf{Who decides when ``nature-inspired'' becomes
``nature-militarised''?} See disagreements below.
\item \textbf{Can ecological security be operationalised without
securitising nature?} NATO's annual assessment and UNEP's nature-based
solutions both assume institutions can measure and act; neither source
resolves how to do so without turning ecosystems into military assets.
\end{itemize}

\subsection{Disagreements and gaps}

\textbf{Climate determinism versus threat multiplier.} The IPCC/NATO
consensus (threat multiplier, weak average effect at present) is
contested by more deterministic readings in popular and some policy
literature, and by a 2026 GMACCC paper framing climate change as ``the
primary threat to global peace and security'' --- a framing the IPCC's
own assessment does not support in that strong form. A paper must not
present determinism as the IPCC position.

\textbf{Militarised biomimicry.} Matthews and Matthews treat mimicry
as continuous with deception and imitation in war; critical scholarship
(for example, Foster and Clapper, ``Swim, Swarm, Strike'', conference
paper \cite{foster2026}) argues biomimetic drones recast nature as a
terrain of permanent counterinsurgency. This is a normative
disagreement, not an empirical one, and a unified programme must address
it rather than assume biomimicry is neutral.

\textbf{Tamper-evidence versus tamper-prevention.} SLSA, Sovereign
Proof, and STAMP all insist tamper-evidence is not prevention; a system
that presents anchored receipts as \emph{preventing} tampering
over-claims. The disagreement in the literature is about how much
deterrence evidence actually provides when the attacker controls the
signer.

\textbf{Evidence quality is uneven across strands.} Sections 3 and 5
have peer-reviewed, replicated results; Section 2's swarm and immune
claims rest heavily on reviews, DARPA programme reports, and preprints;
Section 4's strongest claims are institutional assessments, not
experiments. A unified paper should grade evidence by strand, not blend
them.

\textbf{Gaps found in this survey (not in the sources):} (a) no
peer-reviewed source on active noise cancellation for sonar was located
--- only patent and trade material; (b) no source directly studies a
unified ``nature, security and defense'' programme --- the synthesis in
this section is the author's, built from the strand sources; (c) no
source was found that tests bio-inspired patterns (swarm, immune) as
\emph{software supply-chain} controls --- all platform transfers in
Sections \ref{sec:biomimicry}--\ref{sec:ecological} are analogies and
are labelled as such.

\section{Development Improvements at a Glance}
\label{sec:at-a-glance}

Table \ref{tab:glance} summarises the strongest platform-relevant
takeaway from each strand and its status. Only the verification strand
(Section \ref{sec:verification}) yields improvements stated directly by
its sources; the others are analogies, and are labelled as such.

\begin{table}[h]
\centering
\begin{tabular}{lll}
\toprule
Strand & Strongest platform-relevant takeaway & Status \\
\midrule
1. Biomimicry & Complementary, fail-differently checks; & Analogy, grounded \\
& no single point of failure; quorum / & in BatDeck + OFFSET \\
& multi-party approval & \\
2. Acoustic / sonar & Two-layer noise handling (deterministic & Analogy / \\
& shield + statistical denoising); & Speculative \\
& synthetic-data testing of the verifier; & \\
& hierarchical (cheap local to expensive & \\
& anchored) checks & \\
3. Ecological security & Threat-multiplier risk framing; annual & Analogy \\
& impact assessment; address root-cause & \\
& drivers; minimise collateral / & \\
& data minimisation & \\
4. Verification & SLSA L2/L3 control-plane provenance; & \textbf{Direct} \\
& key must not be reachable by build steps; & \\
& anchor Merkle roots, not raw hashes; & \\
& verify on a schedule; fail closed; & \\
& anchor does not prove authorship; & \\
& passkeys for human signing; dedicated & \\
& service key for the gateway & \\
5. Connection & Deterministic, reproducible trust & \textbf{Direct} \\
& scoring; evidence stays home, only a & (Sovereign Proof \\
& 32-byte commitment leaves; & pattern) \\
& tamper-evidence over tamper-prevention & \\
\bottomrule
\end{tabular}
\caption{Platform improvements at a glance. ``Direct'' means the source
states the pattern and it maps onto the platform's existing controls;
``Analogy'' means the pattern is in the source and the transfer to
software is an inference; ``Speculative'' means a plausible transfer
with no source support for the software application.}
\label{tab:glance}
\end{table}

\section{Limitations and What Could Not Be Verified}
\label{sec:limitations}

\begin{itemize}
\item \textbf{Active noise cancellation (sonar):} no peer-reviewed
source located; only patent and trade summaries. Do not cite ANC as an
established peer-reviewed sonar technique without further research.
\item \textbf{Michigan echolocation paper details:} authors, volume and
pages, and DOI for Patil et al.\ \emph{Journal of Sound and Vibration}
119229 and Kwon et al.\ \emph{IEEE Access} (2025) are taken from the
University of Michigan ARC project listing (index); the papers
themselves were not opened. Verify before final bibliography.
\item \textbf{MDPI bio-inspired blockchain authors:} the abstract page
as opened did not display author names in the extracted text; do not
add authors without re-checking.
\item \textbf{NIST SP 800-193 and Sigstore/in-toto specifics:} read via
secondary summaries (index). The SLSA specification itself was verified
live and should be the primary citation for provenance; Sigstore and
Rekor details should be verified against sigstore.dev and in-toto.io
before final use.
\item \textbf{No source was found for a unified ``nature, security and
defense'' programme as such.} Section \ref{sec:connection}'s synthesis
is an original synthesis from the strand sources, not a sourced claim.
\item \textbf{Platform transfers in Sections
\ref{sec:biomimicry}--\ref{sec:ecological} are analogies.} Only
Section \ref{sec:verification}'s improvements are directly stated by
their sources. The labels must not be removed in the final paper.
\item \textbf{STAMP authors and year:} not captured from the manuscript
PDF as opened; cite by title only until verified.
\item \textbf{Passkeys (FIDO2/WebAuthn):} the WebAuthn Level 3
Recommendation and NIST SP 800-63B-4 were verified live
\cite{webauthn3,nist80063b4}; the FIDO Alliance overview is an index
source \cite{fido-passkeys}. The hybrid proposal
in Section \ref{sec:verification} --- passkeys for human signing, a
dedicated service key for the gateway --- is the author's design
proposal for the E-CO platform, not a sourced claim, and has not been
tested as a software supply-chain control.
\end{itemize}

\begin{thebibliography}{99}

\bibitem{yamada2026} Yamada, Y., Hasegawa, K., Watabe, Y., Sasaki, S.,
Asada, T. \& Hiryu, S. (2026). ``A thermophone based FM sonar for
airborne ultrasonic sensing under drone flight noise.''
\emph{Scientific Reports} 16:26651. DOI: 10.1038/s41598-026-65909-3.
\url{https://www.nature.com/articles/s41598-026-65909-3.pdf} ---
verified live.

\bibitem{saranga2026} Velmurugan, M., Brush, P., Balfour, C., Przybyla,
R.~J. \& Sanket, N.~J. (2026). ``Milliwatt ultrasound for navigation in
visually degraded environments on palm-sized aerial robots.''
\emph{Science Robotics} 11(112), eadz9609. DOI:
10.1126/scirobotics.adz9609.
\url{https://arxiv.org/abs/2603.24699v1} --- verified live (arXiv
abstract).

\bibitem{batdeck2024} M\"uller, H., Kartsch, V., Magno, M. \& Benini,
L. (2024). ``BatDeck --- Ultra Low-power Ultrasonic Ego-velocity
Estimation and Obstacle Avoidance on Nano-drones.''
arXiv:2412.10048v1.
\url{https://arxiv.org/html/2412.10048v1} --- verified live. Earlier
version: arXiv:2403.16696v1.

\bibitem{michigan2025} Patil, G.~U., Kwon, H.~S., Epureanu, B.~I. \&
Popa, B.~I. (2025). ``Synthetically-trained neural networks for shape
classification from measured acoustic scattering.'' \emph{Journal of
Sound and Vibration}, 119229. Kwon, H.~S., Patil, G.~U., Epureanu,
B.~I. \& Popa, B.~I. (2025). ``Specialized Convolutional Neural Network
Models for Echolocation-Based Perception.'' \emph{IEEE Access}. As
listed at
\url{https://arc.engin.umich.edu/research/projects/perception-in-complex-scenes-using-automatically-labeled-sonar-imaging-data-in-synthetic-environments/}
--- index.

\bibitem{wpi2025} Worcester Polytechnic Institute / Sanket, N.~J. NSF
project on bat-inspired echolocation drones (2025--). Reported at
\url{https://i-hls.com/archives/132067} and
\url{https://fixedd.co.uk/bats-inspire-wpi-researchers-to-develop-drones-using-echolocation/}
--- index. Hierarchical reinforcement learning plus physics-informed
deep learning; under 100~mm / 100~g.

\bibitem{matthews2024} Matthews, R. \& Matthews, T. (2024). ``Military
mimicry: the art of concealment, deception, and imitation.''
\emph{Defence and Security Analysis}. DOI:
10.1080/14751798.2024.2352271.
\url{https://research.birmingham.ac.uk/en/publications/military-mimicry-the-art-of-concealment-deception-and-imitation/}
--- index.

\bibitem{igi2025} ``Innovation of Bio-Inspired Materials for
Next-Generation Defense Applications.'' (2025). IGI Global, ch.~5. DOI:
10.4018/979-8-3373-0933-0.ch005.
\url{https://www.igi-global.com/viewtitlesample.aspx?id=382059&ptid=361704&t=Innovation+of+Bio-Inspired+Materials+for+Next-Generation+Defense+Applications}
--- index. Nacre, mantis shrimp, cephalopod camouflage, self-healing.

\bibitem{armadillo2024} ``Biomimetic design and impact simulation of
Al2O3/Al composite armor based on armadillo shell.'' \emph{Scientific
Reports} (2024).
\url{https://www.nature.com/articles/s41598-024-71255-z} --- index.
Armadillo carapace plus triply periodic minimal surfaces; simulation
only, ballistic performance not yet experimentally proven for the TPMS
variant.

\bibitem{darpa2020} DARPA (2020). ``OFFSET Swarm Systems Integrators
Demonstrate Tactics to Conduct Urban Raid.''
\url{https://www.darpa.mil/news/2020/offset-swarm-urban-raid} --- index.
250-system swarms, urban raid field experiment, Camp Shelby.

\bibitem{kim2007} Kim, J., Bentley, P.~J., Aickelin, U., Greensmith,
J., Tedesco, G. \& Twycross, J. (2007). ``Immune system approaches to
intrusion detection --- a review.'' \emph{Natural Computing} 6:413--466.
DOI: 10.1007/s11047-006-9026-4.
\url{https://staff.fmi.uvt.ro/~daniela.zaharie/ma2017/projects/applications/intrusion_detection_systems/AIS%2BIDS/AIS_intrusion_detection.pdf}
--- index.

\bibitem{bala2025} Bala, A. \& El Bhiri, B. (2025). ``A Comparative
Analysis of Immunity-Inspired Cybersecurity Approaches.''
\emph{Engineering Proceedings} 112, 73. DOI:
10.3390/engproc2025112073.
\url{https://mdpi-res.com/d_attachment/engproc/engproc-112-00073/article_deploy/engproc-112-00073.pdf?version=1764318708}
--- index.

\bibitem{swarmmesh2026} Rolle, J.~S. (2026). ``Swarm-Mesh: A
Decentralized Bio-Inspired System for AI Containment.'' Zenodo preprint,
DOI: 10.5281/zenodo.21200799.
\url{https://zenodo.org/records/21200799/files/SwarmMesh_AI_Containment_Preprint_2026-07-04.pdf}
--- index; independent preprint, not peer-reviewed. Quorum-sensing,
resilient consensus, no single point of failure.

\bibitem{mdpi-bioblockchain} ``Bioinspired Blockchain Framework for
Secure and Scalable Wireless Sensor Network Integration in Fog--Cloud
Ecosystems.'' (2024). \emph{Sensors} / MDPI, 14(1):3.
\url{https://www.mdpi.com/2073-431X/14/1/3} --- verified live
(abstract). Authors not captured; do not cite author names. AIS plus
Proof of Adaptive Immunity Consensus; claimed 95.2\% anomaly detection,
500 nodes.

\bibitem{ipcc2022} IPCC (2022). \emph{Climate Change 2022: Impacts,
Adaptation and Vulnerability.} Contribution of Working Group II to the
Sixth Assessment Report. \url{https://www.ipcc.ch/report/ar6/wg2/} ---
index (cited via secondary summaries).

\bibitem{newsecuritybeat} New Security Beat. ``What Next for Climate
Security? Implications From IPCC Working Group II 6th Assessment
Report.''
\url{https://WWW.NEWSECURITYBEAT.ORG/?p=119047} --- index. Also:
\emph{Climatic Change} (2025) 178:125, DOI: 10.1007/s10584-025-03970-z,
\url{https://cdn.cloud.prio.org/files/fdd4e7ad-ad28-4f4b-acc2-ad7d1073606b/s10584-025-03970-z.pdf?inline=true}
--- scientific uptake in UNSC debates 2007--2022; AR6 average climate
effect on conflict judged weak at present, more prominent with
additional warming.

\bibitem{nato2021} NATO (2021). ``NATO Climate Change and Security
Action Plan.'' 14 June 2021.
\url{https://www.nato.int/en/about-us/official-texts-and-resources/official-texts/2021/06/14/nato-climate-change-and-security-action-plan}
--- verified live.

\bibitem{nato-topic} NATO. ``Environment, climate change and
security.''
\url{https://nato.int/en/what-we-do/wider-activities/environment-climate-change-and-security}
--- verified live.

\bibitem{unep2009} United Nations Environment Programme (2009).
\emph{From Conflict to Peacebuilding: The Role of Natural Resources and
the Environment.} ISBN 978-92-807-2957-3.
\url{https://preview-www.cbd.int/financial/doc/unep-conflict.pdf} ---
index.

\bibitem{unep2024} United Nations Environment Programme (2024).
\emph{Nature-based Solutions for Peace: Emerging practice and options
for policymakers.} Nairobi. ISBN 978-92-807-4204-6.
\url{https://wedocs.unep.org/bitstreams/f4db3bd7-e326-481d-af02-4ea887583ff4/download}
--- index.

\bibitem{slsa} SLSA. ``Build: Requirements for producing artifacts.''
Specification v1.2.
\url{https://slsa.dev/spec/v1.2/build-requirements} --- verified live.

\bibitem{filigrane2026} Filigrane (2026-07-15). ``Blockchain
Notarization: When Anchoring Hashes Makes Sense.''
\url{https://www.filigrane.com/blockchain-notarization} --- verified
live.

\bibitem{sovereign2026} Bardwell-Scott, M. (2026). ``A Cryptographically
Signed Biography for Physical Hardware: Tamper-evident chain of custody
for the AI and critical-infrastructure supply chain.'' Sovereign Proof
Technical White Paper v1.3, October 2026.
\url{https://sovereignproof.io/sovereign-proof-whitepaper.pdf} ---
verified live; vendor white paper, not peer-reviewed.

\bibitem{nist800193} Regenscheid, A. (NIST) (2018). \emph{Platform
Firmware Resiliency Guidelines.} NIST SP 800-193, Final, May 2018.
Summarised at
\url{https://www.electronicdesign.com/technologies/industrial/article/55367149/infineon-technologies-secure-tomorrows-data-centers-with-platform-firmware-resiliency}
and
\url{https://softwsp.com/guides/nist-sp-800-193-platform-firmware-resiliency/}
--- index.

\bibitem{sigstore-guide} Sigstore / SLSA / in-toto secondary guides
(index, use primary specs for citation):
\url{https://dev.to/kamal_namdeo/slsa-sigstore-and-provenance-complete-practical-guide-1bl5} ;
\url{https://blog.michaelsam94.com/supply-chain-security-slsa-sbom/} ;
\url{https://hld.handbook.academy/curriculum/security-at-scale/supply-chain-security/}
--- index.

\bibitem{stamp} STAMP: ``Synthesis Tamper-evident Attestation and
Molecular Provenance (STAMP): Cryptographic Molecular Barcoding for DNA
Synthesizers.'' Manuscript PDF.
\url{https://framerforms.s3.us-east-1.amazonaws.com/form-uploads/761ee2a5-94a2-4849-b65f-14b43cabbc10-stamp-cryptographic-molecular-barcoding.pdf}
--- index; preprint, authors not captured.

\bibitem{foster2026} Foster, D. \& Clapper, A. ``Swim, Swarm, Strike:
Biomimetic War Machines and the Cyborg Ecologies of Perpetual
Counterinsurgency.'' Conference paper.
\url{https://nomadit.co.uk/conference/pollen2026/paper/94200} --- index.

\bibitem{webauthn3} W3C Web Authentication Working Group (2026). ``Web
Authentication: An API for accessing Public Key Credentials --- Level
3.'' W3C Recommendation, 25 August 2026.
\url{https://www.w3.org/TR/webauthn-3/} --- verified live. Successor to
Level 2 (8 April 2021); defines the API for strong, attested, scoped,
public key-based credentials. FIDO2 is WebAuthn plus CTAP; a passkey is
a discoverable WebAuthn credential.

\bibitem{nist80063b4} Temoshok, D., Fenton, J.~L., Choong, Y.~Y.,
Lefkovitz, N., Regenscheid, A., Galluzzo, R. \& Richer, J.~P. (2025).
``Digital Identity Guidelines: Authentication and Authenticator
Management.'' NIST Special Publication SP 800-63B-4, July 2025. DOI:
10.6028/NIST.SP.800-63B-4.
\url{https://csrc.nist.gov/pubs/sp/800/63/b/4/final} --- verified live.
Phishing resistance requires cryptographic authentication (\S3.2.5,
WebAuthn as the verifier-name-binding example); syncable authenticators
are admitted at AAL2 but not at AAL3, which requires a non-exportable
key.

\bibitem{fido-passkeys} FIDO Alliance. ``Passkeys'' and ``User
Authentication Specifications.''
\url{https://fidoalliance.org} and
\url{https://fidoalliance.org/specifications-overview/} --- index. FIDO
standards use public-key cryptography to provide phishing-resistant
authentication; every passkey is unique and bound to the online service
domain.

\end{thebibliography}

\end{document}
